% SHARD-ID: RC-07-DELIGNE-VIA-GRAPHS
% SHARD-TITLE: Deligne's Proof, Told Through Regular Graphs
% SHARD-SUMMARY: Reads the Riemann hypothesis over finite fields with the Ihara zeta of a regular graph as the running model, separating what is cheap (rationality, duality, the Perron bound) from what is not (the square-root bound itself).
% SHARD-SUMMARY: States Deligne's pointwise theorem in twisted-graph language with the two-loop bouquet counterexample, and records the three ingredients a variety has and a graph lacks: products, slicing, and the sign.
% SHARD-KEYWORDS: Deligne, Ihara zeta, Ramanujan graph, pointwise purity, big monodromy, tensor powers, the sign, abelian obstruction
% SHARD-NOTES: notes/deligne-via-graphs.md
% SHARD-SCRIPTS: none

\section{Deligne's Proof, Told Through Regular Graphs}
\label{sec:deligne}

Expository. Nothing here is new; the point is to isolate what is actually hard in the proof of the Riemann
hypothesis over finite fields \cite{deligne1974weil}, using the Ihara zeta \cite{ihara1966discrete} of a
regular graph as the running model, and to read off which ingredients this project has and which it does
not.

\subsection{One graph, two sides}

Fix a finite connected $(\qs+1)$-regular graph $X$. The objects are those of \cref{sec:definitions}:
non-backtracking walks and prime cycles (\cref{def:nonbacktracking-walk}), the count $\Ncount{\wl}\ge0$ of
closed non-backtracking walks of length $\wl$, the Ihara zeta with its Euler product
(\cref{def:ihara-zeta}), the Hashimoto operator with $\Ncount{\wl} = \Tr\Hash^{\wl}$ and $\zetaX(\uz) =
\det(1-\uz\Hash)^{-1}$ (\cref{def:hashimoto-operator}), and its compression to the adjacency matrix $\Aadj$
by Ihara--Bass (\cref{def:ihara-bass}, \cite{bass1992ihara}). Side A is the orbit side, side B the small
matrix (\cref{def:side-a-side-b}); passing between them is bookkeeping.

Call $\pm1$ and the pair $\edgeev\in\{\qs,1\}$ from the adjacency eigenvalue $\lambda = \qs+1$ the
\emph{trivial} eigenvalues of $\Hash$, everything else nontrivial. Then
\[
  \Ncount{\wl} \;=\; \qs^{\wl} + 1 + \!\!\sum_{\edgeev\ \mathrm{nontrivial}}\!\! \edgeev^{\wl}
  \;+\; (|E|-|V|)\bigl(1+(-1)^{\wl}\bigr),\]
with extra terms if $X$ is bipartite: the leading term is the growth rate of the $(\qs+1)$-regular tree, and
the nontrivial eigenvalues are the error. Since $\edgeev\edgeev' = \qs$, a pair either has both members of
modulus $\sqrt{\qs}$ (when $|\lambda|\le2\sqrt{\qs}$, where they are complex conjugates) or is real with one
member larger. So all nontrivial $|\edgeev| = \sqrt{\qs}$ iff all nontrivial $|\lambda| \le 2\sqrt{\qs}$ iff
the error in $\Ncount{\wl}$ is $O(\qs^{\wl/2})$: the Ramanujan property (\cref{def:ramanujan-graph}) and the
RH analogue (\cref{def:rh-analogue}) are the same statement, that closed walks spread as evenly as they
possibly could.

Two examples from the founding session. $K_4$ has $\qs=2$, adjacency spectrum $3,-1,-1,-1$, hence $\edgeev =
(-1\pm i\sqrt7)/2$ with $|\edgeev|^2 = 2$: Ramanujan, with $\Ncount{3} = 8+1+6\cdot\tfrac52 = 24$, the
oriented pointed triangles. The prism $C_{16}\times K_2$ has an eigenvalue $2\cos(\pi/8)+1 \approx 2.848 >
2\sqrt2$: a real pole inside the circle and an error growing faster than $2^{\wl/2}$. A long thin graph
always fails: a bottleneck gives a real $\edgeev$ near $\qs$, and walks pile up on one side of it.

\subsection{The same statement for a curve, with one sign changed}

A curve $C$ over $\Fq$ has $\Ncount{\wl}$ points over $\Fqn{\wl}$: again a nonnegative count, again a
generating function $\curvezeta(\uz)$ with an Euler product, now over closed points, the orbits of Frobenius
$\Frob$ on the points over the algebraic closure. Side A is identical in form; side B gives a $2g\times2g$
matrix ($g$ the genus) with
\[
\Ncount{\wl} = \qs^{\wl} + 1 - \Tr\Frob^{\wl} = \qs^{\wl} + 1 - \sum_{i=1}^{2g}\alphaw{i}^{\wl}, \qquad
\curvezeta(\uz) = \frac{\det(1-\uz\Frob)}{(1-\uz)(1-\qs\uz)} .
\]

\begin{center}\begin{tabular}{lll}\toprule
 & graph & curve \\
\midrule
count & closed non-backtracking walks & points over $\Fqn{\wl}$ \\
leading term & $\qs^{\wl}+1$ & $\qs^{\wl}+1$ \\
error term & $+\sum\edgeev^{\wl}$ & $-\sum\alphaw{i}^{\wl}$ \\
pairing & $\edgeev\leftrightarrow\qs/\edgeev$ & $\alphaw{i}\leftrightarrow\qs/\alphaw{i}$ \\
interesting eigenvalues are & poles of $\zetaX$ & zeros of $\curvezeta$ \\
RH & $|\edgeev| = \sqrt{\qs}$ & $|\alphaw{i}| = \sqrt{\qs}$ \\
\bottomrule\end{tabular}\end{center}

Everything matches except the sign of the error term, hence whether the interesting eigenvalues are poles or
zeros. For the graph $\zetaX$ is the reciprocal of a polynomial and every eigenvalue is a pole; for the
curve the trivial eigenvalues $1$ and $\qs$ are poles and the interesting ones are zeros. This is the sign
fixed in \cref{sec:conventions}, $\expsum{n} = -\sum_i\alphaw{i}^n$, and the sign of the explicit formula.
For a variety of dimension $n$ the pattern repeats as an alternating sum over $n+1$ blocks, block $i$ of
modulus $\qs^{i/2}$; the one that matters is the middle, $i=n$.

\subsection{What is cheap and what is not}

\begin{center}\begin{tabular}{lll}\toprule
ingredient & graph & variety \\
\midrule
side B exists & $\Hash$ plus Ihara--Bass & Grothendieck's trace formula \\
the pairing & $\edgeev\edgeev' = \qs$ from Ihara--Bass & Poincar\'e duality \\
trivial bound & Perron: $|\edgeev|\le\qs$ & $|\alphaw{i}|\le\qs$ \\
the bound $\sqrt{\qs}$ & \emph{false in general} & Deligne's theorem \\
\bottomrule\end{tabular}\end{center}

The first three rows are cheap. Rationality of the zeta function is the first row and nothing more; the
trace formula is heavy technology but the same move, that the count is the trace of a power of a finite
matrix. The pairing turns the one-sided bound into equality, so the whole problem is one-sided; and the
trivial bound is just convergence of the Euler product for $|\uz| < 1/\qs$. The gap between $\qs$ and
$\sqrt{\qs}$ is the entire content.

The last row is where graph and variety part. There is no theorem that a $(\qs+1)$-regular graph is
Ramanujan; most are not, the prism is a counterexample, and each Ramanujan construction carries its own
proof: arithmetic (Lubotzky--Phillips--Sarnak \cite{lubotzky1988ramanujan}, through RH for a modular
curve), probabilistic (Friedman), or interlacing (Marcus--Spielman--Srivastava). So the right question is
not how one proves the bound but
what varieties have that graphs lack, which makes every one of them Ramanujan. Three things: varieties can
be \emph{multiplied}, \emph{sliced} over a curve, and their interesting eigenvalues enter the count with a
\emph{minus sign}. A graph has none of the three.

\subsection{The engine: power up, dominate, take roots}

The trick the proof is built on, in its purest graph form: why the Ramanujan bound is $2\sqrt{\qs}$. On the
infinite $(\qs+1)$-regular tree let $c_{2k}$ count closed ordinary walks of length $2k$ from a root, so the
spectral radius of the tree's adjacency operator is $\lim_k c_{2k}^{1/2k}$. Such a walk is a Dyck path
(about $4^k$, up to a polynomial factor) together with a fresh neighbour at each of the $k$ outward steps
($\qs$ choices, $\qs+1$ at the root), so $c_{2k} = 4^k\qs^k$ times a bounded factor times a polynomial in
$k$, and $c_{2k}^{1/2k}\to2\sqrt{\qs}$. Three moves: a nonnegative count dominates a power of an eigenvalue,
$c_{2k} = \langle\delta_{\mathrm{root}},\Aadj^{2k}\delta_{\mathrm{root}}\rangle$; the count is known up to a
loss that does not grow exponentially in $k$; and the $2k$-th root as $k\to\infty$ kills every fixed loss.
This computation is the tree side of the Ramanujan comparison. Deligne makes the same three moves twice; the
nontrivial content is entirely in the second.

\subsection{The pointwise theorem, in twisted-graph language}

Deligne's core theorem is not about a variety but about a curve carrying matrix-valued coefficients. A
\emph{twist} on a base graph $B$ assigns to each directed edge $e$ an invertible $d\times d$ matrix
$\hol(e)$ with $\hol(\bari{e}) = \hol(e)^{-1}$; a closed walk $w$ has holonomy $\hol(w)$, the ordered
product along it, and the twisted zeta is the Artin--Ihara $\Lcurve$-function of \cref{def:artin-ihara},
\[
\Lcurve(\uz,\hol) = \prod_{\pcycle}\det\bigl(1-\uz^{\wl(\pcycle)}\hol(\pcycle)\bigr)^{-1} =
\exp\Bigl(\sum_{w}\tfrac{1}{\wl(w)}\Tr\hol(w)\,\uz^{\wl(w)}\Bigr) = \det(1-\uz\Hash_{\hol})^{-1}.
\]
The quantum Ihara zeta of \cref{sec:quantum-ihara} is this with $B$ a bouquet of $\Dk$ loops and $\hol =
\Ad(\Un{i})$ on the $i$-th loop; a quantum expander is a Ramanujan twist of a bouquet. Two kinds of
eigenvalue now live on side B: the \emph{global} ones, the eigenvalues of $\Hash_{\hol}$ and poles of
$\Lcurve$, analogues of the $\edgeev$; and the \emph{pointwise} ones, the eigenvalues of each
$\hol(\pcycle)$. The twist is \emph{pure with constant $\dweight$} when every eigenvalue of every
$\hol(\pcycle)$ has modulus $\dweight^{\wl(\pcycle)}$: a unitary twist is pure with $\dweight=1$, Deligne's
coefficient systems with $\dweight = \qs^{\wt/2}$ for an integer weight $\wt$ (\cref{def:pure-weight}).

Deligne's pointwise theorem (his Theorem 3.2, in this language): if (a) all traces $\Tr\hol(w)$ are real, in
fact rational; (b) $\hol$ preserves an alternating pairing up to the scalar $\dweight^{2\wl(w)}$, so
eigenvalues of $\hol(\pcycle)$ come in pairs $\alphaw{i}$ and $\dweight^{2\wl(\pcycle)}/\alphaw{i}$; and (c)
the holonomies generate a group as large as (b) permits; then $\hol$ is pure with constant $\dweight$.

The note argues in four steps, for a fixed prime cycle $\pcycle$ and eigenvalue $\alphaw{1}$ of
$\hol(\pcycle)$, the target being $|\alphaw{1}| \le \dweight^{\wl(\pcycle)}$ with (b) supplying equality.
\emph{Power up}: replace $\hol$ by $\hol^{\otimes2k}$, so traces become $(\Tr\hol(w))^{2k}$ and
$\alphaw{1}^{2k}$ is an eigenvalue at $\pcycle$. \emph{Dominate}: by (a) those traces are nonnegative, so
every Euler factor of $\Lcurve_{2k} := \Lcurve(\uz,\hol^{\otimes2k})$ has nonnegative coefficients and
constant term $1$, the product dominates any single factor coefficientwise, and the factor at $\pcycle$ has
a pole at $|\uz| = |\alphaw{1}|^{-2k/\wl(\pcycle)}$ which a dominated series cannot outrun. \emph{The known
pole}: suppose the only poles of $\Lcurve_{2k}$ were the trivial ones, at $|\uz| = 1/(\qs\,\dweight^{2k})$;
then
\[
|\alphaw{1}|^{-2k/\wl(\pcycle)} \;\ge\; 1/(\qs\,\dweight^{2k}) \iff |\alphaw{1}| \;\le\;
\dweight^{\wl(\pcycle)}\,\qs^{\wl(\pcycle)/2k}.
\]
\emph{Take roots}: as $k\to\infty$ the loss $\qs^{\wl(\pcycle)/2k}$ vanishes. The fixed loss is the single
factor $\qs$ in the trivial pole, paid once by the base curve however large $k$ is.

A graph breaks the third step, and must, since the theorem is false for graphs. Take $B$ a bouquet of
two loops with $\hol(a) = \left(\begin{smallmatrix}2&1\\1&1\end{smallmatrix}\right)$, $\hol(b) =
\left(\begin{smallmatrix}1&1\\1&2\end{smallmatrix}\right)$: traces are integers, so (a); determinants
are $1$, so (b) with $\dweight=1$; and the two generate a group as large as $\SL(\mathbb Z)$ permits, so
(c). Yet $\hol(a)$ has eigenvalue $(3+\sqrt5)/2 > 1$, so the twist is not pure: $\Lcurve_{2k} =
\det(1-\uz\Hash_{\hol^{\otimes2k}})^{-1}$ has a pole at $|\uz| = ((3+\sqrt5)/2)^{-2k}$, far inside $1/\qs$,
from a nontrivial eigenvalue of the twisted edge operator. For a graph every side-B eigenvalue is a pole, so
nothing forces the poles to the trivial place. What makes the third step true for a curve is one hypothesis
and one structural fact.

\subsection{Big monodromy, and the abelian failure}

The trivial poles of $\Lcurve_{2k}$ come from invariant vectors: a vector in the $d^{2k}$-dimensional
coefficient space fixed by every holonomy. On such a vector the twist acts like no twist, so $\Lcurve_{2k}$
acquires a copy of the untwisted zeta with its pole at $1/\qs$, rescaled by the scalar the twist puts on the
vector; if that scalar is known, so is the pole. The pairing (b) supplies invariants for free, by
contracting the $2k$ tensor slots in pairs, each contraction giving a pole at exactly
$1/(\qs\dweight^{2k})$. Hypothesis (c), \emph{big monodromy}, says these are the only invariants: if the
holonomies fill out the whole group preserving the pairing, invariant theory makes every invariant of
$\hol^{\otimes2k}$ a combination of pairwise contractions, so no other trivial pole appears.

\begin{observation}[The abelian obstruction is inside Deligne's proof]
\label{obs:abelian-obstruction}
\claimstatus{obs:abelian-obstruction}{sketched}
Commuting holonomies produce uncontrolled invariants in tensor powers, exactly as commuting Kraus unitaries
never give an expander: commuting $\hol(w)$ have common eigenvectors, and a common eigenvector with
eigenvalue $\chi(w)$ becomes invariant in $\hol^{\otimes N}$ as soon as $\chi^N = 1$, carrying an unknown
$\chi$-dependent scalar, so the trivial poles sit where the argument cannot control them. Any Ramanujan lift
of the prime dilations must break the commutativity by a quotient, as \cite{lubotzky1988ramanujan} does with
the arithmetic lattice and as Deligne does with the shears of a generic slicing.
\end{observation}

In Deligne's setting the commuting case is a coefficient system with abelian monodromy, which can have
arbitrary pointwise weights; (c) excludes it. In this project it is the observation of the founding session
that commuting Kraus unitaries never give an expander, the adjoint action of an abelian group having as many
invariants as it has joint eigenprojectors: the same obstruction on both sides of the dictionary.
Geometrically, big monodromy is verified through shears. The holonomy around each bad slice is the identity
plus a rank-one nilpotent along one vector, those vectors are carried to one another by the holonomy group,
and a theorem of Kazhdan and Margulis says a group generated by one conjugacy class of shears and acting
without invariant subspaces is automatically as large as the pairing allows. Big monodromy is the statement
that the bad slices are all alike, which holds because they are generic.

\subsection{The sign}

For a graph, side B is the one matrix $\Hash_{\hol}$ and every eigenvalue is a pole. For a curve with
coefficients the trace formula has two blocks of opposite sign,
\[
\sum_{w:\,\wl(w)=\wl}\!\!\Tr\hol(w) = \Tr\bigl(\Frob^{\wl}\,\big|\,\text{invariant}\bigr) -
\Tr\bigl(\Frob^{\wl}\,\big|\,\text{interesting}\bigr), \quad
  \Lcurve(\uz,\hol) = \frac{\det(1-\uz\Frob\mid\text{interesting})}{\det(1-\uz\Frob\mid\text{invariant})}.\]
The invariant block is the one big monodromy computes, and it produces the poles; the interesting block
produces zeros, and zeros do not obstruct convergence, so the radius of convergence of $\Lcurve_{2k}$ is set
by the invariant block alone. That is the third step, and it is the minus sign of the table above:
positivity of counts is available on both sides, but on the graph side it constrains the sum of everything,
which is the Perron bound and no more, while on the curve side the interesting eigenvalues are subtracted,
so positivity bounds them above by the trivial ones and even tensor powers squeeze the bound to equality.
The two-loop bouquet is not a curve because its walk counts cannot be written with the interesting part
subtracted.

\subsection{From pointwise to global: slice, bound, multiply}

\emph{Slice.} Take a one-parameter family of hyperplane sections of a variety $X$ of dimension $n$; the
parameter runs over a curve and all but finitely many slices are smooth of dimension $n-1$. Up to pieces
already understood, the middle block of $X$ is the interesting block of that base curve with coefficients,
the coefficient at a parameter value being the middle block of the slice there. In graph terms this is the
Artin--Ihara factorisation (\cref{def:artin-ihara}) used in \cref{sec:quantum-ihara}, with the fibre
replaced by a variety of one dimension less. The coefficients satisfy (a), (b), (c): the pairing is
Poincar\'e duality on the slice, big monodromy is the shear argument, and the traces are rational because
the zeta of a slice is rational with rational coefficients, off which the blocks known by induction are
peeled. So they are pure with $\dweight = \qs^{(n-1)/2}$.

\emph{Bound.} With pure coefficients the Euler product over the base converges absolutely for $|\uz| <
1/(\qs\dweight)$, there being about $\qs^{\wl}/\wl$ prime cycles of length $\wl$ each with holonomy of norm
$\dweight^{\wl}$. So the interesting block has no eigenvalue larger than $\qs\dweight = \qs^{(n+1)/2}$,
where RH wants $\qs^{n/2}$: a loss of one factor $\sqrt{\qs}$, paid by the base being a curve, independent
of $n$.

\emph{Multiply.} The product $X^k$ has dimension $nk$ and its middle block contains the $k$-fold tensor
products of the middle block of $X$, so the products $\alphaw{1}\cdots\alphaw{k}$ are among its eigenvalues.
The bound applied to $X^k$ still loses one factor $\sqrt{\qs}$, because the base of a slicing is always a
curve, whence
\[
  |\alphaw{i}|^k \le \qs^{(nk+1)/2} \implies
  |\alphaw{i}| \le \qs^{n/2}\,\qs^{1/2k} \xrightarrow{\ k\to\infty\ } \qs^{n/2},\]
with duality giving the reverse inequality: the engine a second time, count, fixed loss, roots. The
pointwise theorem is what makes the last move legitimate, since the slices of $X^k$ have dimension $nk-1$,
far above anything the induction on dimension knows, so their purity must come from rationality and big
monodromy alone. That is why the core theorem is stated for coefficients on a curve: it certifies
coefficients in any dimension at the price of one curve. Even for a curve $C$ the argument passes through
$C^k$ for all $k$; the surface $C\times C$ is where Weil's proof lived, and the products are the mechanism.
The blocks $i\ne n$ follow by induction from lower-dimensional sections and from duality.

\subsection{What the proof does not do}

It never exhibits a Hermitian structure. Weil's proof for curves does: by the Hodge index theorem the
intersection form on $C\times C$ is definite on the relevant part, which makes $\Frob/\sqrt{\qs}$ unitary
for an explicit positive inner product on the middle block and produces a Hilbert--P\'olya operator
(\cref{def:hilbert-polya}). In graph terms, on the nontrivial part $\Hash$ is $\sqrt{\qs}$ times a unitary
in some inner product, which for the block
$\left(\begin{smallmatrix}\lambda&-\qs\\1&0\end{smallmatrix}\right)$ holds exactly when $|\lambda| <
2\sqrt{\qs}$. Grothendieck's plan was to prove a definite form in all dimensions, the standard conjectures,
and deduce RH as Weil had; that plan is still open. Deligne bypassed it: no form, no operator, only
positivity and a limit. RH over finite fields is therefore a theorem, while in dimension above one the inner
product that would make Frobenius normal is not known to exist.

The Artin--Schreier computation of \cref{sec:artin-schreier} sits on the Weil side of this divide:
$\Tmat\Tmat^\dagger = \qs I$ is manifest from Parseval for the additive characters
(\cref{prop:as-unitarity}), an explicit inner product. Deligne's proof is what one has when no such
structure is in sight. Two honesty notes the note records: the Ramanujan property of the LPS graphs of
\cref{sec:weil-lps} needs only the curve case, since Eichler--Shimura transports Hecke eigenvalues to
Frobenius eigenvalues of a modular curve, so Weil's theorem suffices there and Deligne's method is what
extends it to higher weight; and the technology behind ``side B exists'' is the largest single piece of the
proof by page count, though not by idea.

\subsection{What this says for the project}

\begin{observation}[The three ingredients, and no Hermitian form]
\label{obs:deligne-ingredients}
\claimstatus{obs:deligne-ingredients}{sketched}
Deligne's proof uses three ingredients that a Hilbert--P\'olya approach would have to replace: products,
under which eigenvalues multiply while the trivial loss stays fixed; slicing over a curve, which keeps that
loss at one factor $\sqrt{\qs}$ in every dimension; and the sign, which makes the interesting eigenvalues
zeros rather than poles, so that positivity of the counts bounds them. It exhibits no Hermitian form and no
operator.
\end{observation}

\begin{itemize}
  \item \textbf{Which ingredients $\zeta$ has.} The zeros are subtracted from the prime count, so
$\zeta$ is on the curve side of the table, not the graph side, and the positivity $\vM(n)\ge0$ is the right
kind. The sign and the positivity are present; the product is absent, since there is no $X\times X$ with
$\zeta$ as its slice, so the multiply step has nothing to act on. This is the standard diagnosis, stated in
the terms of this section.
  \item \textbf{The number-field shadow of the pointwise theorem is Rankin--Selberg.} For a modular
form the Euler product of $\Lcurve(f\times\bar f)$ has nonnegative coefficients and a known pole, and
dominating one local factor gives $|a_\pp| \le \pp^{(k-1)/2}\cdot\pp^{1/2}$, the same fixed loss; each
symmetric-power $\Lcurve$-function whose poles are known shaves it (Kim--Sarnak's $7/64$ is the current
state). Knowing the poles of all the tensor-power $\Lcurve$-functions is the analogue of big monodromy; over
number fields that is Langlands functoriality, unproved.
  \item \textbf{The abelian obstruction is inside the proof} (\cref{obs:abelian-obstruction}), on
  both sides of the dictionary.
  \item \textbf{Ramanujan without Hermiticity is possible.} Deligne's route and the interlacing
route both reach the spectral bound with no Hilbert--P\'olya operator. For the steering question of where
Hermiticity enters, this is the cleanest evidence that it need not enter at all; what must enter is a
positivity plus an operation under which eigenvalues multiply at fixed loss.
\end{itemize}
